\section{Semantic Search with HNSW}
\marginpar{TD}

Modern applications increasingly require not only exact key-value retrieval but also \textit{semantic search}: retrieving documents that are related in meaning rather than merely matching the same words or keys. This is typically achieved by mapping text or images to embedding vectors in which semantically similar items lie close together. Retrieval then becomes a similarity-search problem: given a query vector, find the $k$ stored vectors that are closest under a chosen distance metric. This is the $k$-Nearest Neighbour ($k$-NN) problem. Although exact search can be solved by a linear scan in $O(n \cdot d)$ time, this quickly becomes too costly for large collections and high-dimensional embeddings. \textit{Approximate Nearest Neighbour} (ANN) methods therefore trade a small amount of recall for much lower query latency.

\subsection{HNSW Overview}
{\textit{Hierarchical Navigable Small World}} (HNSW), introduced by Malkov and Yashunin, organises vectors in a multi-layer proximity graph. The bottom layer contains all nodes, while higher layers contain progressively fewer nodes. This hierarchy supports efficient long-range navigation at the upper levels and local refinement near the bottom. For the full algorithm and its theoretical background, we refer to the original HNSW paper\cite{8594636}.

\medskip
\noindent
Each inserted vector is assigned a random maximum level. In our implementation, layer 0 stores up to $2M$ neighbours per node, whereas higher layers store up to $M$ neighbours. The graph also maintains a single \emph{entry point}, which serves as the starting node for both insertion and search.

\subsubsection{Insertion, Search, and Deletion}
Insertion and search follow the standard two-stage HNSW structure. During insertion, the algorithm starts from the current entry point and greedily descends through the upper layers, repeatedly moving to a closer neighbour. Once it reaches the levels where the new node must be linked, it performs a broader local search to collect candidate neighbours and then applies the standard HNSW pruning heuristic before creating the final edges.

\medskip
\noindent
Search proceeds similarly. It first performs greedy descent through the upper layers and then switches to a beam search of width \texttt{efSearch} at layer 0. This beam search explores several promising candidates in parallel and returns the top-$k$ nearest vectors found, yielding a much lower practical query cost than an exact linear scan.

\medskip
\noindent
Deletion in the current persisted index is not implemented as pure soft deletion. The native in-memory core can mark nodes as deleted, but the public \texttt{HNSWIndex.delete()} path rebuilds the collection-local persisted graph without the removed document and writes the retained nodes back to disk. The removed node and its vector slot are freed immediately, so deleted routing nodes do not remain in the persisted graph after commit. This makes deletion substantially more expensive than insertion or search, but avoids reopen-time connectivity problems that occurred when only soft-deleted nodes were filtered out.

\subsection{Native Execution Model}
Our initial prototype implemented HNSW entirely in TypeScript, but profiling showed that the performance-critical work was dominated by managed-runtime overhead. We therefore moved the graph computation to Rust while retaining persistence in TypeScript. Rust now handles the main computational kernel of the index, including distance evaluation, graph traversal, insertion, and search, whereas TypeScript remains responsible for storage, recovery, commit orchestration, and the current persisted delete path.

\medskip
\noindent
This separation keeps the existing storage architecture intact while moving the most expensive graph operations to native code. After the move to Rust, additional optimisations reduced the remaining host-side overhead, for example by batching neighbour retrieval across layers, reusing cached buffers, and writing neighbour lists directly in the hot path. These changes did not alter the algorithm itself, but reduced the cost of synchronising modified graph state back to persistent storage.

\subsection{Persistent Representation and Crash Consistency}
The HNSW index uses two persistent structures inside the shared \texttt{FreeBlockFile}: a vector heap for raw \texttt{float32} embeddings and graph pages for node topology. In the vector heap, each component occupies 4 bytes, so a vector of dimension $d$ requires exactly $4d$ bytes of payload. Vector pages additionally store a small header and slot bitmap for allocation tracking.

\medskip
\noindent
Graph pages store node metadata and neighbour lists. Each node contains its document identifier, level, vector pointer, deleted flag, and the neighbour references for all stored layers. Neighbour references are encoded compactly as a page identifier plus a slot identifier. The resulting serialized node size is
$$
\text{nodeSize} = 12 + (1 + 2M \cdot 6) + L(1 + M \cdot 6),
$$
where $L$ is the maximum stored level. With $M = 16$ and $L = 6$, as used in the wrapper-level HNSW benchmarks, one serialized node occupies 787 bytes.

\medskip
\noindent
Crash consistency is provided by the shared WAL-backed \texttt{FreeBlockFile}. Vector pages, graph pages, records, and ordinary index pages all belong to the same atomic commit domain, so a single \texttt{commit()} persists semantic and non-semantic state together. Insertion still follows a logical \emph{data before index} order, but the fundamental persistence guarantee comes from the shared atomic commit mechanism.

\subsection{Integration with SimpleDBMS}
HNSW is integrated per collection through an optional \texttt{HnswConfig}, which specifies the embedding dimension, a text embedding function, and optionally an image embedding function. When a document is inserted, the DBMS determines whether semantic indexing should use image content or text content. If a supported image field is present and an image embedding function is configured, the first such image is embedded; otherwise, the DBMS embeds the concatenation of the public string fields. The resulting vector is then inserted into the collection's HNSW index.

\medskip
\noindent
At query time, the public API exposes \texttt{semanticSearch} for text queries and, when an image embedding function is configured, \texttt{semanticImageSearch} for image queries. In both cases, the query is embedded first, then matched against the collection's HNSW index, after which the matching document locations are resolved back to DBMS records and returned through \texttt{SimpleDbms} as JSON-serialised documents. Because the semantic configuration itself is not persisted, the caller must pass the same configuration again on \texttt{open()} if semantic search should remain available after reopening.

\subsection{Performance}
We evaluated the implementation using four benchmark families: a wrapper-level HNSW complexity benchmark through \texttt{HNSWIndex}, a DBMS-level insert-overhead benchmark comparing B\textsuperscript{+}-tree-only inserts against B\textsuperscript{+}-tree+HNSW inserts inside \texttt{SimpleDbms}, recall-latency sweeps on SIFT1M, GIST1M, and a synthetic random 1M 128D dataset, and a direct insert-only comparison between the Rust-backed implementation and the legacy TypeScript version. All runs used real storage with 4096-byte blocks. The wrapper-level HNSW, recall, and Rust-versus-TypeScript benchmarks used $M = 16$, \texttt{efConstruction} $= 200$, \texttt{efSearch} $= 128$, and \texttt{maxStoredLevel} $= 6$. The DBMS insert-overhead benchmark used the same block size and HNSW $M$ and \texttt{efConstruction}, but a deterministic mock 384-dimensional encoder and the DBMS defaults for the remaining HNSW parameters because it measures insert cost rather than search quality.
SIFT1M and GIST1M are standard Texmex ANN benchmark datasets with public ground truth \cite{jegou2011product}. SIFT1M contains one million 128-dimensional image descriptors and serves as the baseline workload, while GIST1M contains one million 960-dimensional descriptors and is included to test the index on a harder high-dimensional case.
\\
\\
\texttt{All benchmarks were run on an Intel 12th Gen Core i7-1260P laptop processor}

\paragraph{Wrapper-level HNSW complexity.}
The first benchmark exercised the persisted \texttt{HNSWIndex} directly on SIFT1M, recording checkpointed insert and search latencies every 1,000 insertions. Over the full build, insert and search operations scale gracefully. Initial inserts and queries complete in well under $0.1$\,ms. By one million vectors, these times rise logarithmically to roughly $1.0$\,ms for inserts and $0.5$\,ms for searches. As shown in Figure~\ref{fig:hnsw-complexity-overview}, both insertion and search costs increase gradually with database size and remain in the sub-millisecond to low-millisecond range even at one million vectors.

\begin{figure}[H]
    \centering
    \includegraphics[width=0.92\linewidth]{figures/hnsw/complexity_overview.png}
    \caption{Wrapper-level HNSW complexity on SIFT1M through \texttt{HNSWIndex}. The four panels show insert time and search time both as functions of database size and normalised by $\log_2 N$.}
    \label{fig:hnsw-complexity-overview}
\end{figure}

\paragraph{DBMS insert overhead.}
For the DBMS-level benchmark we compare ordinary document inserts with only the B\textsuperscript{+}-tree structures enabled against inserts that additionally maintain a collection-local HNSW index inside \texttt{SimpleDbms}. To isolate index-maintenance cost from model inference, this benchmark uses a deterministic mock 384-dimensional encoder instead of an external embedding model. The B\textsuperscript{+}-tree-only insertions are highly efficient, completing in fractions of a millisecond regardless of collection size. Adding the HNSW index increases the cost from roughly $0.3$\,ms initially to around $4.1$\,ms at one million documents. Because the pure B\textsuperscript{+}-tree cost is extremely small, the relative overhead factor increases from roughly $6\times$ initially to more than $150\times$. As shown in Figure~\ref{fig:dbms-insert-overhead-ratio}, the structural cost of maintaining the semantic index heavily dominates the plain document-insert path once the collection becomes large.

\begin{figure}[H]
    \centering
    \includegraphics[width=0.92\linewidth]{figures/dbms/insert_overhead_ratio.png}
    \caption{Insert-overhead ratio inside \texttt{SimpleDbms}, shown as the ratio between B\textsuperscript{+}-tree+HNSW insert time and B\textsuperscript{+}-tree-only insert time as the database grows.}
    \label{fig:dbms-insert-overhead-ratio}
\end{figure}
\paragraph{Embedding model latency.}
The insert-overhead benchmark above intentionally uses a deterministic mock encoder that returns a pre-computed random vector in under $0.01$\,ms per call, isolating pure index-maintenance cost from model inference. In real deployments the embedding call is a significant additional cost per insert. Table~\ref{tab:embedding-speed} shows representative per-call encoding latencies for a selection of models, measured on the test machine over 150 samples with a 5-call warm-up using the \texttt{@huggingface/transformers} ONNX runtime. Even the fastest real model (MiniLM-L3 q8 at $1.84$\,ms mean) adds latency comparable to the raw HNSW insert cost at low collection sizes, while larger base-class models such as nomic-embed-text add more than an order of magnitude relative to the mock. The overhead ratios reported in the DBMS benchmark therefore represent a lower bound; in production the per-document cost would be dominated by the encoder. All tested models are freely available on the Hugging Face Hub under permissive open-source licenses (such as Apache 2.0 or MIT).

\begin{table}[H]
  \centering
  \caption{Embedding model latency on the test machine (Intel i7-1260P, 150 samples, 5 warm-up calls each). Sorted by mean latency.}
  \label{tab:embedding-speed}
  \begin{tabular}{lcr}
    \toprule
    Model & dim & mean (ms) \\
    \midrule
    \textit{mock-random-384d}           & 384 & 0.01  \\
    paraphrase-MiniLM-L3-v2 (q8)        & 384 & 1.84  \\
    all-MiniLM-L6-v2 (q8)               & 384 & 2.43  \\
    all-MiniLM-L6-v2 (fp32)             & 384 & 3.26  \\
    bge-small-en-v1.5 (q8)              & 384 & 6.10  \\
    nomic-embed-text-v1.5 (q8)          & 768 & 16.12 \\
    \bottomrule
  \end{tabular}
\end{table}
\paragraph{Recall-latency trade-off.}
Recall benchmarks were run on SIFT1M, GIST1M, and a synthetic random dataset of one million 128-dimensional vectors. The results are shown in Figure~\ref{fig:hnsw-recall-combined}. On SIFT1M, recall@10 (the fraction of true nearest neighbours successfully found in the top 10 results) rose from $0.96206$ at \texttt{efSearch} $= 64$ to $0.99898$ at \texttt{efSearch} $= 500$, with mean latency increasing from $0.289$ ms to $1.749$ ms. On GIST1M, which is higher-dimensional and more difficult, recall@10 increased from $0.7526$ to $0.9635$ over the same range, with mean latency rising from $1.500$ ms to $8.649$ ms. By contrast, the random 1M 128D dataset remained difficult even at large beam widths: recall@10 increased only from $0.072$ to $0.290$. This is consistent with the fact that HNSW benefits from meaningful geometric structure in the data, whereas fully random vectors offer little neighbourhood structure to exploit.

\begin{figure}[H]
    \centering
    \includegraphics[width=0.82\linewidth]{figures/hnsw/recall_vs_ef_search_combined.png}
    \caption{Recall@10 as a function of \texttt{efSearch} for SIFT1M, GIST1M, and random 1M 128D. Structured datasets benefit much more strongly from larger beam widths than the random baseline.}
    \label{fig:hnsw-recall-combined}
\end{figure}

\paragraph{Rust versus legacy TypeScript.}
Finally, we compared the current Rust-backed implementation with the legacy TypeScript HNSW implementation on the same SIFT1M base set, using identical storage settings and HNSW parameters. The benchmark measured only insert operations inside the main loop; the reported insert mean therefore excludes the final commit, while the wall-clock time includes it. In the current 1,000,000-vector run, the Rust implementation achieved a mean insert time of $0.813$ ms and a wall-clock build time of $844.13$ s, whereas the legacy TypeScript implementation required $3.144$ ms per insert on average and $3169.16$ s wall-clock time. This corresponds to a practical insertion speedup of approximately $3.87\times$. As shown in Figure~\ref{fig:hnsw-insert-compare}, the Rust-backed implementation remains consistently faster throughout the build. This comparison should be read as an implementation-level benchmark rather than as a claim that both paths are identical in every internal detail, but it does show that moving the computational kernel to Rust yields a substantial real-world speedup within the same overall storage architecture.

\begin{figure}[H]
    \centering
    \includegraphics[width=0.9\linewidth]{figures/hnsw/insert_compare_sift.png}
    \caption{Insert-only comparison on SIFT1M between the current Rust-backed implementation and the legacy TypeScript implementation. The Rust-backed implementation stays consistently faster throughout the build.}
    \label{fig:hnsw-insert-compare}
\end{figure}
